\documentclass{handout}\usepackage{handoutSetup}\usepackage{handoutShortcuts}  

\begin{document}

\handoutHeader

\begin{verbatimwrite}{\jobname.title}
The Ramsey/Cass-Koopmans (RCK) Model with Government
\end{verbatimwrite}

\centerline{\LARGE \input \jobname.title}
\medskip 
\centerline{\url{http://www.econ2.jhu.edu/people/ccarroll/public/lecturenotes}}

\medskip \medskip

This handout solves the Ramsey/Cass-Koopmans (RCK) model with government.\footnote{The treatment is similar
to that in \cite{blanchard&fischer:text}; see that source for more details.}  For simplicity we assume
no technological progress, population growth, or depreciation, and a 
continuum of consumers with mass 1 distributed on the unit interval
as per \handoutF{Aggregation}.

Consider first the case where government is financed by constant 
lump-sum taxes of amount $\Tax$ per period, and spending is at rate 
$\mathbf{x}$ per period, and suppose the government has a balanced budget 
requirement so that $\mathbf{x} = \Tax$ in every period and there is no 
government debt.  We suppose further that government spending yields 
no utility.  The individual's optimization problem (leaving out the $i$ 
subscripts that we used in the previous handout, but understanding that 
they are implicitly present) is now
\begin{equation}\begin{gathered}\begin{aligned}
        \int_{0}^{\infty} \util(c_{t}) e^{-\timeRate t}
\end{aligned}\end{gathered}\end{equation}
subject to the household DBC
\begin{equation}\begin{gathered}\begin{aligned}
        \dot{a}_{t} & =  r_{t} a_{t} + w_{t} - c_{t} - \Tax, \label{eq:hhbc}
\end{aligned}\end{gathered}\end{equation}
which has Hamiltonian representation
\begin{displaymath}
        \Ham(a_{t},c_{t},\lambda_{t}) = \util(c_{t}) + (r_{t}a_{t} + w_{t} - c _{t} - \Tax)\lambda_{t}
\end{displaymath}

The first Hamiltonian optimization condition requires $\partial 
\Ham_{t}/\partial c_{t} = 0$:
\begin{equation}\begin{gathered}\begin{aligned}
        c_{t}^{-\rho} & =  \lambda _{t}  \label{eq:H1A} \\
        -\rho c^{-\rho-1}_{t}\dot{c}_{t} & =  \dot{\lambda}_{t}
\end{aligned}\end{gathered}\end{equation}

The second Hamiltonian optimization condition requires:
\begin{equation}\begin{gathered}\begin{aligned}
        \dot{\lambda}_{t} & =  \timeRate\lambda_{t} - (\partial \Ham_{t}/\partial a_{t})\lambda_{t}  \\
         & =  \timeRate\lambda_{t} - \lambda_{t} r_{t}
\\      \dot{\lambda}_{t}/{\lambda_{t}} & =  (\timeRate-r_{t})
\\  \dot{c}_{t}/c_{t}  & =  \rho^{-1}(r_{t}-\timeRate).
\end{aligned}\end{gathered}\end{equation}

Finally, the household's behavior must satisfy a transversality 
constraint, which is equivalent to the intertemporal budget constraint:
\begin{equation}\begin{gathered}\begin{aligned}
        \int_{0}^{\infty} c_{t}\mathfrak{R}^{-1}_{t} & =  a_{0}+\int_{0}^{\infty} \mathbf{y}_{t} \mathfrak{R}^{-1}_{t}-\int_{0}^{\infty} \Tax \mathfrak{R}^{-1}_{t}
\\  \mathbf{C}_{0} & =        a_{0}+\mathbf{Y}_{0}-\mathbf{T}_{0}
\end{aligned}\end{gathered}\end{equation}
which says that the present discounted value of consumption must equal 
the current net physical wealth plus human wealth minus the PDV of 
taxes.

Now consider the problem from the standpoint of a social planner who has the same utility function as the 
individual consumers.  If the social planner wants to spend a constant
amount $\mathbf{x}$ per period, the social planner's budget constraint is
\begin{equation}\begin{gathered}\begin{aligned}
        \dot{k}_{t} & =  \fFunc(k_{t}) - \depr k_{t} - c_{t} - \mathbf{x}
\end{aligned}\end{gathered}\end{equation}
which reflects the fact that the social planner divides total net output
between consumption and government spending.  This leads to Hamiltonian
\begin{equation}\begin{gathered}\begin{aligned}
        \Ham(k_{t},c_{t},\lambda_{t})  & =  \util(c_{t}) + \lambda_{t}(\fFunc(k_{t}) - \depr k_{t} - c_{t} - \mathbf{x}),
\end{aligned}\end{gathered}\end{equation}
yielding the first order condition 
\begin{equation}\begin{gathered}\begin{aligned}
        \dot{c}_{t}/c_{t} & =  \rho^{-1}(\fFunc(k_{t})-\depr-\timeRate).
\end{aligned}\end{gathered}\end{equation}

Now recall from the handout on decentralizing the RCK model
that
\begin{equation}\begin{gathered}\begin{aligned}
        \fFunc(k_{t}) & =  \hat{r}_t k_{t} + w_{t}
\end{aligned}\end{gathered}\end{equation}
where the gross return on capital $\hat{r}_t$ is equal to the net
return $r_t$ plus the depreciation rate.  

Thus, the social planner's DBC is:
\begin{equation}\begin{gathered}\begin{aligned}
        \dot{k}_{t} & =  \fFunc(k_{t})-\depr k_{t} - c_{t} - \mathbf{x}
\\          & =  \hat{r}_{t} k_{t}     - \depr k_{t} + w_t - c_t - \mathbf{x}
\\          & =  r_{t} k_{t} + w_{t} - c_{t} - \mathbf{x},
\end{aligned}\end{gathered}\end{equation}
which is equivalent to the household's budget constraint 
\eqref{eq:hhbc} when $a_{t}=k_{t}$ and $\Tax=\mathbf{x}$.  As discussed 
in \handoutG{DecentralizingRCK}, $a_t=k_{t}$ must hold in equilibrium for identical households, 
and $\Tax = \mathbf{x}$ was the balanced budget assumption that we started 
off with.

Note that the $\dot{c}_t=0$ locus in the phase diagram is unchanged by 
changing $\Tax$ and $\mathbf{x}$.  However, the $\dot{k}_{t}=0$ locus is 
shifted down by amount $\Tax = \mathbf{x}$.  

Now what happens if the government does not face a balanced budget
requirement?  Specifically, suppose we continue to have the same
constant amount of spending per period but now want to consider the
effect of allowing taxes to vary over time, which we denote by a
subscript on $\Tax_{t}$.  Suppose $d$ is the level of government bonds
(debt); the government's Dynamic Budget Constraint is
\begin{equation}\begin{gathered}\begin{aligned}
        \dot{d}_{t} & =  \mathbf{x}+r_{t} {d}_{t}-\Tax_{t},
\end{aligned}\end{gathered}\end{equation}
which says that debt must rise by the amount by which spending exceeds 
taxes.

The government's IBC will be the integral of its DBC:
\begin{equation}\begin{gathered}\begin{aligned}
        {d}_{0}+\int_{0}^{\infty}\mathbf{x} \mathfrak{R}^{-1}_{t}   & =  \int_{0}^{\infty} \Tax \mathfrak{R}_{t}^{-1}   \\
        {d}_{0}+\mathbf{X}_{0} & =  \mathbf{T}_{0}
\end{aligned}\end{gathered}\end{equation}
and we assume ${d}_{0}=0$ so that the government starts out with no debt 
(to maintain comparability with the previous example).

The DBC of the idiosyncratic family also changes.  They can now own 
either capital $k_{t}$ or government debt ${d}_{t}$.  If the family is to be 
indifferent between the two forms of assets, the interest rate must be 
the same.
\begin{equation}\begin{gathered}\begin{aligned}
        c_{t} + \dot{a}_{t} & =  w_{t} + r_{t} a_{t} - \Tax_{t} \\
        a_{t} & =  k_{t} + {d}_{t}.
\end{aligned}\end{gathered}\end{equation}

Now the family's IBC becomes
\begin{equation}\begin{gathered}\begin{aligned}
        \int_{0}^{\infty} c_{t} \mathfrak{R}^{-1}_{t} & =  k_{0}+{d}_{0}+\mathbf{Y}_{0}-\int_{0}^{\infty} \Tax  _{t} \mathfrak{R}^{-1}_{t}
\\      & =  k_{0}+{d}_{0}+\mathbf{Y}_{0}-\mathbf{T}_{0} \label{eq:famibc}.
\end{aligned}\end{gathered}\end{equation}

Note: Nowhere in this equation does the time path of taxes matter;
all that matters is the PDV of taxes.  And the time path of taxes
also does not enter the $\dot{c}/c$ equation.  Thus, the path 
of consumption over time is unaffected by the path of taxes over time!

This is not so surprising when you realize that it is simply the 
Ricardian equivalence proposition in this perfect foresight framework.

However, now consider the case where there is a tax on capital income
at rate $\tau$.  Furthermore, for simplicity suppose that the
government rebates all of the tax revenue in a lump sum per capita,
and suppose depreciation $\depr=0$.  Thus the household budget
constraint becomes
\begin{equation}\begin{gathered}\begin{aligned}
        \dot{a}_t & =  r_{t}(1-\tau)a_{t}+w_{t}-c_{t}+\mathbf{z}_{t}
\end{aligned}\end{gathered}\end{equation}
where $\mathbf{z}_{t}$ is the per-capita size of the lump-sum rebates,
\begin{equation}\begin{gathered}\begin{aligned}
        \mathbf{z}_{t} & =  \tau r_{t} k_{t}.
\end{aligned}\end{gathered}\end{equation}

The household's Hamiltonian becomes
\begin{equation}\begin{gathered}\begin{aligned}
        \Ham_{t}(c_{t},a_{t},\lambda_{t}) & =  \util(c_{t})+\lambda_{t}(r_{t}(1-\tau)a_{t}+w_{t}-c_{t}+\mathbf{z}_{t}).
\end{aligned}\end{gathered}\end{equation}

The crucial difference between this situation and the previous one is
that now the effective rate of return on saving has been decreased, so
that $\partial \Ham_{t}/\partial a_{t}$ is now $r_{t}(1-\tau)$ rather
than $r_{t}$.  Ultimately this produces a consumption Euler equation
of
\begin{equation}\begin{gathered}\begin{aligned}
        \dot{c}_{t}/c_{t} & =  \rho^{-1}((1-\tau)r_{t}-\timeRate)
\end{aligned}\end{gathered}\end{equation}
which implies that the economy will be in equilibrium at
\begin{equation}\begin{gathered}\begin{aligned}
        \fFunc(\bar{k})(1-\tau) & =  \timeRate  \\
        \fFunc(\bar{k}) & =  \timeRate
\end{aligned}\end{gathered}\end{equation}
so that the equilibrium level of the marginal product of capital is higher,
and the capital stock must therefore be lower, than before the capital
taxation was instituted.  

Notice, however, that because the taxes are being rebated, the {\it aggregate}
budget constraint does not change when the tax is imposed:
\begin{equation}\begin{gathered}\begin{aligned}
        \dot{k}_t & =  r_{t}(1-\tau)k_{t}+w_{t}-c_{t}+\tau k_{t}r_{t}  \\
         & =  r_{t}k_{t}+w_{t}-c_{t}.
\end{aligned}\end{gathered}\end{equation}

Thus the social planner will choose exactly the same amount of
consumption as before the tax was instituted.

The crucial point is that if an individual household saves more and 
thus causes next year's capital stock to be a bit higher, the {\it 
personal} benefit to that household is essentially zero.  The higher 
taxes that the household will pay next year will be distributed to the 
entire population in a lump sum, so the saver will get nothing.  The 
higher saving of this individual household is basically a positive 
externality from the point of view of the other consumers in the 
economy.  However, if the social planner forces the economy as a whole 
to save more, the social planner receives all of the extra tax 
revenue.

 
\input handoutBibMake
\end{document}

Similarly, the social planner's IBC is
\begin{equation}\begin{gathered}\begin{aligned}
        \int_{0}^{\infty} c_{t} \mathfrak{R}^{-1}_{t} & =  k_{0}+\mathbf{Y}_{0}-\int_{0}^{\infty} \mathbf{x}   \\
\\       & =  k_{0}+\mathbf{Y}_{0}-\mathbf{X}_{0}
\end{aligned}\end{gathered}\end{equation}
but since $k_{0}=a_{0}$ and since taxes equal spending in every period 
$\mathbf{T}_{0}=\mathbf{X}_{0}$ so that the social planner's IBC is identical to the 
household's IBC.

